\section{An exact weight-ten commutator normalization}
\label{sec:weight-ten-normalization}

Use $G=D_9H$ and $J_9=D_9\cap H$. For $g=(1,f)\in D_9$, define
\[
 m_{10}(g)=[S^5TST^3]f.
\]
This mixed coefficient detects the next even-degree bracket.

\begin{theorem}\label{thm:weight-ten-exact-image}
The full and arithmetic images are
\[
 m_{10}(D_9)=m_{10}(J_9)=\mathbf Z_2.
\]
For $D_{10}=D_9\cap\ker m_{10}$ and $J_{10}=D_{10}\cap H$,
\[
 D_9=D_{10}J_9,\qquad G=D_{10}H,\qquad
 B\simeq D_{10}/\overline{[G,D_{10}]J_{10}}.
\]
The subgroup $D_{10}$ is proper, closed and normal, and its words
start in degree at least eleven.
\end{theorem}

\subsection{The primitive leading bracket}

Let $L_w$ denote the standard Lyndon Lie brackets in $S<T$, and set
\[
 \psi_3=L_{SST}-L_{STT},\qquad
 \mathcal L_{10}=\{\psi_3,\mathcal L_7\},
\]
where $\mathcal L_7$ is the primitive degree-seven vector displayed
above. As before,
\[
 \{P,Q\}=[P,Q]+\partial_P(Q)-\partial_Q(P),\qquad
 \partial_P(S)=0,\quad\partial_P(T)=[T,P].
\]
All coefficients of $\mathcal L_{10}$ are integral. Two exact
coefficient calculations are
\begin{equation}\label{eq:weight-ten-unit-coefficient}
 [L_{SSSSSTSTTT}]\mathcal L_{10}=-141,
 \qquad [S^5TST^3]\mathcal L_{10}=-141.
\end{equation}
In particular it is primitive over $\mathbf Z_2$.

\begin{lemma}\label{lem:weight-ten-leading-line}
For every multiplier-one GT word starting in degree ten, the
homogeneous term of degree ten is $A\mathcal L_{10}$ with
$A\in\mathbf Z_2$. It vanishes exactly when $m_{10}$ vanishes.
\end{lemma}

\begin{proof}
The following explicit finite certificate proves the necessary
rational line. Collect the $99+186=285$ standard Lyndon group
commutators of weights ten and eleven. Through degree eleven
write
\[
 F=1+\sum_{10\le|w|\le11}h_w(C_w-1).
\]
Products of two nonconstant factors start in degree twenty, so
this expression is exact at the indicated precision. Define the
cycle forms
\[
 I_v=[v]\bigl(F(Y,X)+F(X,Y)-2\bigr),\qquad
 II_v=[v]\bigl(F(Z,X)+F(Y,Z)+F(X,Y)-3\bigr),
\]
where $X=1+S$, $Y=1+T$ and $Z=Y^{-1}X^{-1}$. Encode the rows by
$2^{|v|}+\operatorname{bin}(v)$ and
$4096+2^{|v|}+\operatorname{bin}(v)$, respectively.

Adjoin one leading Artin pentagon row $P$, with code $8192$.
It is the coefficient of the output word
\[
 11123234324
\]
on $X_1$. Explicitly, use the infinitesimal derivations
$d_{ij}(X_i)=[X_i,X_j]$, $d_{ij}(X_j)=[X_j,X_i]$, with the other
letters fixed. Evaluate the degree-ten Lie polynomial on the
five pairs
\[
 (d_{12},d_{23}+d_{24},+),\quad
 (d_{13}+d_{23},d_{34},+),\quad
 (d_{23},d_{34},-),\quad
 (d_{12}+d_{13},d_{24}+d_{34},-),\quad
 (d_{12},d_{23},-),
\]
and take the indicated coefficient of the signed sum. This is a
necessary coefficient of the full pentagon: a word starting in
degree ten contributes first in Artin output degree eleven, so
only the leading derivations enter. The common sign relating the
ordinary pure Artin action to the displayed $d_{ij}$ is raised to
the even power ten. The row is zero on all degree-eleven columns.

Let $A_R$ be the following $222$ rows, in the written order:
\[
\begin{aligned}
R=(&1025,1027,1029,1031,1033,1035,1037,1039,1041,1043,1045,1047,1049,1051,\\
&1053,1059,1061,1063,1065,1067,1069,1073,1075,1077,1079,1081,1083,1085,\\
&1093,1095,1097,1099,1101,1107,1109,1113,1115,1117,1125,1129,1131,1133,\\
&1139,1171,1173,1197,1203,1237,2049,2051,2053,2055,2057,2059,2061,2063,\\
&2065,2067,2069,2071,2073,2075,2077,2079,2081,2083,2085,2087,2089,2091,\\
&2093,2095,2097,2099,2101,2103,2105,2107,2109,2115,2117,2119,2121,2123,\\
&2125,2127,2129,2131,2133,2135,2137,2139,2141,2147,2149,2151,2153,2155,\\
&2157,2161,2163,2165,2167,2169,2185,2187,2189,2195,2197,2199,2201,2203,\\
&2205,2211,2213,2215,2217,2219,2221,2227,2229,2233,2247,2249,2251,2253,\\
&2259,2261,2265,2267,2277,2281,2341,2347,2349,2355,2357,2379,2381,2387,\\
&2389,5121,5127,5131,5133,5139,5141,5145,5155,5157,5161,5167,5169,5189,\\
&5193,5207,5211,5291,5301,6145,6147,6149,6151,6153,6155,6157,6159,6161,\\
&6163,6165,6167,6169,6171,6177,6179,6181,6183,6185,6187,6189,6191,6193,\\
&6195,6197,6199,6211,6213,6215,6217,6219,6221,6225,6227,6229,6231,6233,\\
&6235,6237,6245,6249,6251,6253,6281,6283,6285,6291,6293,6295,6297,6299,\\
&6307,6309,6313,6315,6317,6325,6345,6347,6357,6437,6443,8192).
\end{aligned}
\]
Let $C$ be the first $222$ columns. Write $\ell_j$ for the
$99$ degree-ten Lyndon coordinates of the explicitly defined
bracket $\mathcal L_{10}$. Its twelfth coordinate is $-141$.
Define the $98$ target rows, of length $285$, by
\[
 T_j=e_j-\frac{\ell_j}{\ell_{12}}e_{12}
       \qquad(1\le j\le99,\ j\ne12).
\]
The exact finite identities are
\[
 \det A_{R,C}=-63256059346747392\equiv37\pmod{101},
 \qquad T_C A_{R,C}^{-1}A_R=T.
\]
Every entry is specified by standard group commutators, geometric
inversion through degree eleven, and the displayed derivation
recursion. Thus the identity forces
$h_j=(\ell_j/\ell_{12})h_{12}$ for every GT word. Its rational
leading scalar is $A=h_{12}/(-141)$, which belongs to
$\mathbf Z_2$ because $141$ is odd. Equation
\eqref{eq:weight-ten-unit-coefficient} gives $m_{10}=-141A$.
This proves the assertion without inferring a lift from the
necessary equations.
\end{proof}

\subsection{An arithmetic unit}

\begin{lemma}\label{lem:weight-ten-arithmetic-unit}
There is an arithmetic commutator $w_{10}\in J_9$ such that
\[
 m_{10}(w_{10})=-141\cdot729\,m_3(\sigma)\tau_7(\sigma)
                                  \in\mathbf Z_2^\times.
\]
Moreover $w_{10}=[\sigma,k]$ for some $k\in D\cap H$.
\end{lemma}

\begin{proof}
Use the established arithmetic element $\sigma\in H\cap M$ with
$\rho_2(\sigma)=2$ and every odd-prime Kummer coordinate zero.
Both $a=m_3(\sigma)$ and $\tau_7(\sigma)$ are odd. Choose
$h\in H$ with $\lambda_2(h)=3$, set $t=\lambda(h)$, and put
\[
 r_3=h\sigma h^{-1}\sigma^{-27},\qquad
 r_5=h r_3h^{-1}r_3^{-243}.
\]
These are multiplier-one arithmetic elements. The first has
$c_3=m_3=0$, so its parameter starts in degree five. The second
has $c_3=c_5=0$, so its parameter starts in degree seven. These
vanishings use the cycle calculations before imposing any radical
normalization.

The weight-seven beta law gives
\[
 c_7(r_5)=\frac{(3^7-3^5)(3^7-3^3)}{360}\tau_7(\sigma)
        =16\cdot729\tau_7(\sigma).
\]
For a word starting in degree seven, $c_7=m_7$. The degree-seven
leading relation, which holds for every multiplier-one GT word
starting in that degree, therefore gives
\[
 \operatorname{in}_7(f_{r_5})=729\tau_7(\sigma)\mathcal L_7.
\]
No normalized weight-seven divisibility claim is being applied to
$r_5$. Its parameter is allowed to have a primitive leading term.
The leading parameter of $\sigma$ is $a\psi_3$. Hence the
commutator
\[
 w_{10}=[\sigma,r_5]
\]
starts in degree ten and has leading term
\[
 729a\tau_7(\sigma)\{\psi_3,\mathcal L_7\}
       =729a\tau_7(\sigma)\mathcal L_{10}.
\]
Every Kummer homomorphism on $M$ vanishes on this commutator. Its
lower Magnus coefficients also vanish, so $w_{10}\in J_9$.
The unit assertion follows from
\eqref{eq:weight-ten-unit-coefficient}.

For the last assertion, the all-prime conjugation law gives
\[
 \rho_q(r_5)=(t-243)(t-27)\rho_q(\sigma)
\]
for every rational prime $q$. Put $c=(t-243)(t-27)$ and
$k=r_5\sigma^{-c}$. The continuous profinite power is defined for
this $c\in\widehat{\mathbf Z}$, and all Kummer coordinates of $k$
are zero. Thus $k\in D\cap H$. Since powers of $\sigma$ commute
with $\sigma$, one has $[\sigma,k]=[\sigma,r_5]=w_{10}$.
\end{proof}

\begin{proof}[Proof of Theorem~\ref{thm:weight-ten-exact-image}]
The leading coefficient $m_{10}$ is an additive continuous
homomorphism on $D_9$, and conjugation acts by
$\lambda_2^{10}$. The arithmetic image contains the unit furnished
by the preceding lemma. Being compact, it is all of
$\mathbf Z_2$, as is the full image. Equality gives
$D_9=D_{10}J_9$ and $G=D_{10}H$; the usual compact gluing argument
identifies the relative quotient $B$ as displayed. The scalar law
proves normality and the unit proves properness. The leading-line
lemma shows that the words of $D_{10}$ start in degree at least
eleven.
\end{proof}

\subsection{The tenth character}

\begin{proposition}\label{prop:eta-ten-extension-kernels}
The invariant character $\eta_{10}=m_{10}\pmod2$ on $D_9$ extends
invariantly to $D_8$ by
\[
 \alpha_{10}(1,f)=[S^5TST^3]f\pmod2,
\]
but has no invariant extension to $D$ or $N$. The original kernel is
\[
 \ker\!\left(H^1(N,\mathbf F_2)^G
                  \longrightarrow H^1(D_{10},\mathbf F_2)\right)
 =\left(\bigoplus_{q\text{ odd prime}}\mathbf F_2\psi_q\right)
                  \oplus\mathbf F_2\eta_3.
\]
The joint raw-coefficient map $(m_9,m_{10})$ on $D_8$ has image
$512\mathbf Z_2\times\mathbf Z_2$ already on $J_8$, with kernel
$D_{10}$. Its invariant quadratic-character kernel on restriction
to $D_{10}$ is $\mathbf F_2\eta_9\oplus\mathbf F_2\alpha_{10}$.
\end{proposition}
\begin{proof}
Indeed the exact weight-nine theorem puts every degree-nine leading
vector on $D_8$ in $32\mathbf Z_2\mathcal L_9$. Modulo two, such
words therefore start in degree at least ten, making the displayed
coefficient additive and invariant under the odd multiplier.

In the other direction, $\eta_{10}$ has no invariant extension to
$D$ or $N$. The arithmetic witness is $w_{10}=[\sigma,k]$ with
$k\in D$, so every such invariant character kills it, whereas
$\eta_{10}(w_{10})=1$. Consequently the original character kernel
is still
\[
 \ker\!\left(H^1(N,\mathbf F_2)^G
                  \longrightarrow H^1(D_{10},\mathbf F_2)\right)
 =\left(\bigoplus_{q\text{ odd prime}}\mathbf F_2\psi_q\right)
                  \oplus\mathbf F_2\eta_3.
\]
To see the last assertion, restrict a character in this kernel to
$D_9$. It is zero or $\eta_{10}$, since $D_9/D_{10}\simeq\mathbf Z_2$;
the second case cannot extend to $N$. The previously proved kernel
on $D_9$ then applies.

Finally, extend $m_{10}$ to $D_8$ by the same raw coefficient.
The map $(m_9,m_{10})$ is additive there, since all
nonconstant product and substitution corrections start above
degree ten. The arithmetic witnesses $w_9,w_{10}$ show that its
image is $512\mathbf Z_2\times\mathbf Z_2$, with kernel $D_{10}$.
Thus the invariant quadratic-character kernel on $D_8$ is
$\mathbf F_2\eta_9\oplus\mathbf F_2\alpha_{10}$.

\end{proof}

\subsection{A group-theoretic extension test in weight ten}
\label{subsec:eta-ten-extension-threshold}

\begin{proposition}\label{prop:eta-ten-no-extension-d7}
The character $\eta_{10}$ on $D_9$ extends invariantly to $D_8$
by $\alpha_{10}$, but has no $G$-invariant extension to $D_7$.
Consequently it has no such extension to $D_5$, $D_3$, $D$ or $N$.
\end{proposition}

\begin{proof}
The extension to $D_8$ was proved above. We give an obstruction
using only class-two commutator identities and the previously
established leading terms.

For $r\ge3$, let $\mathcal F^r$ be the subgroup of
$M=\ker\lambda$ whose pro-$2$ Magnus parameter starts in degree
at least $r$. The filtered normalized group law gives
\[
 [\mathcal F^r,\mathcal F^s]\subseteq\mathcal F^{r+s}.
\]
One has $M=\mathcal F^3$ and
$[M,M]\subseteq D_7\subseteq\mathcal F^8$. Thus
$[M,[M,M]]\subseteq\mathcal F^{11}$, and the quotient
$M/\mathcal F^{11}$ has nilpotence class at most two.
Moreover,
\[
 [\mathcal F^5,\mathcal F^5]\subseteq\mathcal F^{11}.
\]
Indeed its possible degree-ten leading term is the Ihara bracket
of two multiples of $\psi_5$, and this is zero. Vanishing over
$\mathbf Q_2$ implies vanishing in the integral leading module,
which is torsion-free.

Keep the arithmetic $\sigma,h$ of the weight-ten construction, and
put
\[
 a_0=\sigma,\qquad a_1=h\sigma h^{-1},\qquad
 a_2=h^2\sigma h^{-2},\qquad c_{ij}=[a_i,a_j].
\]
The weight-three law shows that
\[
 r=a_1a_0^{-27},\qquad q=a_2a_0^{-729}
\]
belong to $\mathcal F^5$. Therefore $[r,q]=1$ in
$M/\mathcal F^{11}$. All commutators are central in this quotient,
and expanding this identity gives
\[
 1=[a_1a_0^{-27},a_2a_0^{-729}]
   =c_{12}c_{01}^{729}c_{02}^{-27}.
\]
Hence $c_{12}=c_{02}^{27}c_{01}^{-729}$ there.

The arithmetic weight-eight and weight-ten witnesses are
\[
 w_8=c_{01},\qquad
 w_{10}=[a_0,(a_2a_1^{-27})(a_1a_0^{-27})^{-243}].
\]
The same class-two calculation gives
\[
 w_{10}=c_{02}c_{01}^{-270}\pmod{\mathcal F^{11}}.
\]
Consequently the word
\[
 z=h w_8h^{-1}w_8^{-6561}=c_{12}c_{01}^{-6561}
\]
satisfies
\[
 z=c_{02}^{27}c_{01}^{-7290}
   =w_{10}^{27}\pmod{\mathcal F^{11}}.
\]
Since $z\in J_7$ and its parameter agrees through degree ten with
that of $w_{10}^{27}$, one has $z\in J_9$ and
\[
 m_{10}(z)=27m_{10}(w_{10}),\qquad\eta_{10}(z)=1.
\]
On the other hand, every $G$-invariant character
$\chi:D_7\to\mathbf F_2$ satisfies
\[
 \chi(z)=\chi(h w_8h^{-1})-6561\chi(w_8)=0,
\]
because $6561$ is odd. Thus such a character cannot restrict to
$\eta_{10}$ on $D_9$. Any extension to one of the larger listed
subgroups would restrict to an extension on $D_7$, so those are
also impossible.
\end{proof}

\begin{corollary}\label{cor:weight-ten-character-kernels}
The invariant quadratic-character kernels on restriction to
$D_{10}$ are
\[
\begin{array}{c|l}
 \text{domain}&\text{kernel}\\ \hline
 D_9&\mathbf F_2\eta_{10}\\
 D_8&\mathbf F_2\eta_9\oplus\mathbf F_2\alpha_{10}\\
 D_7&\mathbf F_2\eta_8\oplus\mathbf F_2\eta_9\\
 D_5&\mathbf F_2\eta_7\oplus\mathbf F_2\alpha_8\\
 D_3&\mathbf F_2\eta_5\\
 N&\displaystyle\left(\bigoplus_{q\text{ odd prime}}\mathbf F_2\psi_q\right)
                    \oplus\mathbf F_2\eta_3.
\end{array}
\]
\end{corollary}

\begin{proof}
The first two rows follow from the exact scalar and joint maps.
For any larger listed domain, a character killed on $D_{10}$
restricts on $D_9$ to zero or $\eta_{10}$. The latter is impossible
by the preceding proposition. Such a character therefore already
kills $D_9$, and the previously proved weight-nine kernel table
applies.
\end{proof}

